% ==========================================================================
% When Search Becomes Structure
%
% Thorsten Fuchs & Anna-Sophie Fuchs
% Panta Rhei Research Program
% May 2026
% ==========================================================================
\documentclass{pantarhei-researchnote}

\usetikzlibrary{positioning,arrows.meta,calc,fit,shapes.geometric}

\input{tau-macros}

\hypersetup{
  pdftitle={When Search Becomes Structure},
  pdfauthor={Thorsten Fuchs and Anna-Sophie Fuchs},
  pdfsubject={Guarded bridge note between Wolfram's ruliological P-vs-NP microcosms and τ-admissible witness construction},
  pdfkeywords={Panta Rhei, Category τ, ruliology, P vs NP, τ-admissibility, τ-admissible witness construction, witness search, address construction, CRT reassembly, computational complexity, Stephen Wolfram, Ruliad}
}

% Paper-specific macros
\newcommand{\Wolfram}{Wolfram}
\newcommand{\Ruliad}{\ensuremath{\mathcal{R}_{\mathrm{uliad}}}}
\newcommand{\TTM}{\ensuremath{\mathrm{TTM}}}
\newcommand{\TATM}{\ensuremath{\mathrm{TATM}}}
\newcommand{\Adm}{\ensuremath{\mathrm{Adm}_{\tau}}}
\newcommand{\Padm}{\ensuremath{\tau\text{-}\mathrm{P}_{\mathrm{adm}}}}
\newcommand{\NPadm}{\ensuremath{\tau\text{-}\mathrm{NP}_{\mathrm{adm}}}}
\newcommand{\TauP}{\ensuremath{\tau\text{-}\mathrm{P}}}
\newcommand{\TauNP}{\ensuremath{\tau\text{-}\mathrm{NP}}}
\newcommand{\Ezero}{\ensuremath{E_0}}
\newcommand{\Eone}{\ensuremath{E_1}}
\newcommand{\Etwo}{\ensuremath{E_2}}
\newcommand{\Ethree}{\ensuremath{E_3}}
\newcommand{\CRT}{\ensuremath{\mathrm{CRT}}}
\newcommand{\Prim}{\ensuremath{\operatorname{Prim}}}
\newcommand{\ZFC}{\ensuremath{\mathrm{ZFC}}}
\newcommand{\ClassP}{\ensuremath{\mathrm{P}}}
\newcommand{\ClassNP}{\ensuremath{\mathrm{NP}}}
\newcommand{\CategoryTau}{Category~\(\tau\)}

\begin{document}

\lane[pr-lane-verify]{Panta Rhei Research \textperiodcentered\ Computation and Readout}
\title{When Search Becomes Structure}
\subtitle{\texorpdfstring{A guarded bridge note between Wolfram's ruliological P-vs-NP microcosms and $\tau$-admissible witness construction}{A guarded bridge note between Wolfram's ruliological P-vs-NP microcosms and tau-admissible witness construction}}
\noteedition{RC1 \textperiodcentered\ May 2026}

\author{%
  Thorsten~Fuchs
  \and
  Anna-Sophie~Fuchs%
}
\date{May 2026}

\footerseries{Research Note}
\footerslug{search-becomes-structure-categorical}
\footerdate{May 2026}

\begin{abstract}
\noindent
Stephen Wolfram's 2026 essay on P vs NP and ruliological computation
studies finite computational microcosms: small Turing-machine classes,
runtime outliers, isolate machines, nondeterministic speedups,
everything-machine reachability, and computational irreducibility
\cite{Wolfram2026RuliologicalPNP}. This Research Note reads that essay as
an external bridge surface for a bounded Panta Rhei question: when does
search stop being blind search and become structure? The answer here is
not a slogan about the classical Clay problem. It is the scoped
τ-admissibility claim of the computation/admissibility spine:
under finite interface width, witness search is treated as address
construction rather than blind enumeration \cite{PR-III}. The central
comparison is asymmetric. Wolfram makes computational difficulty visible
from below; the τ-framework specifies when witness search becomes
construction from within. The note does not claim that Wolfram supplies
validation for \CategoryTau, that the ruliad is \CategoryTau, or that the
classical \ZFC/Turing-machine P-vs-NP problem is resolved.
\end{abstract}

\keywords{P vs NP \textperiodcentered\ ruliology \textperiodcentered\
computational irreducibility \textperiodcentered\ witness search
\textperiodcentered\ address resolution \textperiodcentered\
computational bi-square \textperiodcentered\ tau-admissibility
\textperiodcentered\ Category $\tau$}

\maketitle

% =======================================================================
\section{Introduction: the bridge and the boundary}
\label{sec:intro}
% =======================================================================

The P-vs-NP problem is famously easy to state and famously hard to locate.
In its orthodox form it asks whether every decision problem whose positive
instances can be verified in polynomial time can also be solved in
polynomial time by a deterministic Turing machine. Yet that tidy statement
leaves many structural choices in the background: what counts as a
machine, what counts as a representation, what kind of nondeterministic
choice is being granted, how observers inspect computations, and what
kind of infinity is being used when asymptotic limits are taken.

\Wolfram's ruliological essay is valuable because it refuses to leave all
of that invisible. It does not try to settle the full question. Instead it
asks what can be learned by explicitly enumerating finite computational
microcosms: small one-sided Turing machines, the integer functions they
compute, the runtime profiles they exhibit, the hard cases that appear
inside fixed rule classes, and the ways nondeterministic branching changes
reachable behavior. This is not orthodox proof theory, but it is not mere
illustration either. It is empirical theoretical computer science in the
literal sense: finite rule universes are sampled, measured, and compared.
The value of Wolfram's ruliological framing is that it makes
computational difficulty visible without first forcing it into a single
orthodox proof template.

The τ-framework reaches the same danger from another direction. It asks
what has to be true before search can be called search at all. In the
framework's computation/admissibility spine, P-vs-NP-style questions are
native only at an operational layer where there are agents, steps, and an
efficiency measure. Once that typing is in place, the decisive distinction
is not simply deterministic versus nondeterministic. It is admissible
versus inadmissible. A τ-admissible verifier has finite interface width:
there is a finite primorial depth beyond which the computation sees
nothing new. Under that condition, a witness is no longer a flat bitstring
certificate. It is a canonical address whose per-prime components can be
resolved and reassembled by the Chinese remainder structure. The slogan
``search becomes structure'' means exactly this: search becomes
construction only where a finite address architecture has been earned.

\prcallout[CLAIM BOUNDARY]{%
This is a guarded bridge note, not a proof note. Wolfram's essay is used
as a public comparator for finite computational difficulty. The formal
collapse introduced below is internal to τ-admissible witness problems,
after finite interface width and address structure have been defined.
This is an internal equality for τ-admissible witness problems. It does
not settle orthodox P vs NP over unrestricted Turing-machine / ZFC
encodings. The note does not treat the ruliad and \CategoryTau\ as the
same object.}

The claim tiers are deliberately separated. Wolfram supplies the external
comparator. The τ-admissible collapse is an internal framework claim. Any
τ-native or physical-computation extension would add further assumptions.
Any orthodox Clay/ZFC export would require a separate encoding, size
measure, reduction, and proof-system bridge.

\begin{widetable}[t]
\caption{Claim-tier map for the note. The tiers are separated near the
front because the same words can otherwise sound much stronger than the
manuscript intends.}
\label{tab:claim-tiers}
\prtablefont
\begin{tabularx}{\textwidth}{p{0.20\textwidth}Y Y}
\toprule
\textbf{Tier} & \textbf{What this note may say} & \textbf{What it must not imply} \\
\midrule
External Wolfram comparator & Wolfram's ruliological essay is a
valuable public anchor for finite computational microcosms, runtime
wildness, nondeterministic reach, and the ruliad/everything-machine
limit \cite{Wolfram2026RuliologicalPNP,WolframInstituteRuliology}. &
It must not imply that Wolfram supplies validation for \CategoryTau,
proves any τ theorem, or resolves the classical P-vs-NP problem. \\
Internal τ claim & Inside the τ computation/admissibility spine, finite
interface width and CRT witness decomposition support a scoped
admissibility-collapse theorem \cite{PR-III}. & It must not imply that
the internal equality is the orthodox
\(\ClassP=\ClassNP\) theorem. \\
Orthodox P-vs-NP boundary & Any export to ordinary Turing-machine
languages, bit-length polynomial time, reductions, and \ZFC\ proof is a
separate bridge problem. & It must not imply that the Clay problem is
solved by this note. \\
Future export bridge & Later work may compare finite ruliological
benchmarks with τ-interface-width diagnostics and explicit witness
decompositions. & It must not imply that benchmark analogies count as
validation, proof, or a universal solution recipe. \\
\bottomrule
\end{tabularx}
\end{widetable}

The contribution of this note is a disciplined comparison. Section
\ref{sec:wolfram} reconstructs the Wolfram anchor in its own terms.
Sections \ref{sec:tau-layer}--\ref{sec:bisquare-collapse} reconstruct the
τ computation source spine: computation layer, τ-Tower Machine, interface
width, witness search, computational bi-square, and admissibility
collapse. Sections \ref{sec:finite-interface}--\ref{sec:ruliad} compare
the two frameworks along four axes: finite exploration, nondeterminism,
irreducibility, and maximal reachability. Section \ref{sec:guardrails}
then makes the orthodox boundary explicit. The closing section turns the
comparison into a research program rather than a headline.

\subsection{Anchor selection and reading method}

Wolfram's essay is an unusually good anchor for this note because it has
the right kind of distance from the τ-framework. It is close enough to
share the core intuition that rule space, observers, and finite
computational microcosms matter. It is distant enough that it cannot be
mistaken for an internal source, a friendly paraphrase, or a validation
surface. The essay is not written in the language of \(\tau\)-addresses,
primorial towers, CRT decomposition, or product-meet diagrams. It is
therefore useful as an external comparator: it presses on the same nerve
without using the same vocabulary.

The selection method is conservative. A paper or essay is a good anchor
for this note only if it has three properties. First, it must discuss
P-vs-NP difficulty through actual computational structure, not only
through slogans about hard problems. Second, it must distinguish finite
experiments from the full orthodox problem, so that the bridge does not
begin with an overclaim. Third, it must expose a genuine contrast between
search, reachability, and construction. Wolfram's essay satisfies all
three. It builds finite machine microcosms; it repeatedly keeps the
classical problem open; and it shows that nondeterministic reach can be
studied without yet giving a deterministic construction method.

The reading method is correspondingly narrow. When this note says that
Wolfram makes a difficulty visible, it means that the essay supplies a
finite observable analogue: a runtime profile, a machine-class relative
outlier, a multiway speedup, or an everything-machine limit. When this
note says that the τ spine answers a different question, it means that
the internal source base imposes a typed admissibility criterion before
claiming collapse. The comparison is not a two-column proof by analogy.
It is a map of where two neighboring frameworks ask similar questions and
where their answers stop being the same.

\subsection{Three separations used throughout}

To keep the paper readable, the following separations are used
throughout.

\begin{enumerate}[leftmargin=1.2em,itemsep=1pt,topsep=2pt]
\item \emph{Reach is not construction.} A short path in a multiway graph
is not yet a deterministic procedure for finding that path.
\item \emph{Finite enumeration is not finite interface width.} A bounded
machine class is an experimental microcosm; a finite-width verifier is an
internal admissibility condition.
\item \emph{Internal collapse is not orthodox export.}
\(\Padm=\NPadm\) belongs to the τ-admissible fragment. The classical
\(\ClassP\) versus \(\ClassNP\) problem over Turing machines and
\ZFC-style proof remains outside the claim of this note.
\end{enumerate}

\begin{widetable}[t]
\caption{The external anchor reconstructed as finite ruliological
microcosms. The right column records the claim discipline used throughout
this note.}
\label{tab:wolfram-microcosms}
\prtablefont
\begin{tabularx}{\textwidth}{p{0.17\textwidth}Y Y Y}
\toprule
\textbf{Surface} & \textbf{What Wolfram studies} & \textbf{Use in this note} & \textbf{Guardrail} \\
\midrule
Finite machines & One-sided Turing machines with fixed state/color
classes computing integer functions. & A concrete rule universe whose
function and runtime structure can be inspected. & Not the same model as
τ-Tower Machines. \\
Runtime outliers & Functions whose runtime profiles jump or contain
dramatic hard inputs. & A finite view of why asymptotic behavior can be
hard to infer. & Outliers are not universal lower bounds. \\
Isolate machines & Machines that are the unique or class-isolated
computers of a function within a fixed finite class. & Bounded evidence
for non-shortcut behavior. & Isolation is not τ-inadmissibility. \\
Nondeterminism & Multiway choice can shorten some computations, sometimes
dramatically, but not uniformly. & A comparison surface for branch
availability. & A branch is not yet deterministic witness construction. \\
Everything machine & A limiting finite-class object that includes all
rule cases and reaches outputs by minimal paths. & A contrast for maximal
nondeterministic reach. & Maximal reach is not τ construction. \\
Ruliad limit & The ruliad is the entangled limit of all possible
computations, sampled by bounded observers. & A philosophical comparator
for observer-bounded readout. & The ruliad is not \CategoryTau. \\
\bottomrule
\end{tabularx}
\end{widetable}

% =======================================================================
\section{Wolfram's finite ruliological exploration}
\label{sec:wolfram}
% =======================================================================

The first obligation in a bridge note is fidelity. Wolfram's essay is not
an appendix to our framework. It has its own problem, vocabulary, and
method. The essay begins from the observation that P vs NP, in full form,
is an infinite question: it quantifies over all possible inputs and all
possible Turing machines. Wolfram's move is to ask finite subquestions
inside that infinite problem. What happens if one fixes a small machine
class and explicitly enumerates its behavior? What functions appear? Which
machines compute them? What runtime profiles occur? When does adding
states, colors, or nondeterministic alternatives change the answer? These
are the finite subquestions Wolfram uses to replace the global problem
with inspectable microcosms \cite{Wolfram2026RuliologicalPNP}.

\subsection{Finite machines as empirical theory}

The machines in the essay are deliberately concrete. Inputs and outputs
are read as integers; machines are specified by small state/color classes;
halting and undefined behavior are tracked rather than swept away. This
turns a famous asymptotic question into a family of finite maps:
machine-class, computed function, runtime profile, and fastest-known
representative. Within a fixed class, one can establish class-relative
facts that are hard to obtain in ordinary complexity theory. For example,
one can ask whether any machine of a given size computes a particular
function faster than another, or whether a function is represented by only
one machine in that finite class \cite{Wolfram2026RuliologicalPNP}.

The word ``relative'' matters. A lower bound inside a fixed small machine
class is not a lower bound for all Turing machines. Wolfram says this
explicitly in practice, because larger machines can sometimes compute the
same finite mapping by a lookup-style or more specialized rule. The result
is not a naive claim that finite hard cases are hard forever. It is a
careful demonstration that hardness is sensitive to the representation and
rule universe in which it is measured \cite{Wolfram2026RuliologicalPNP}.

For the present note, that class relativity is not a defect. It is the
point. A finite computational universe makes visible what a purely
asymptotic statement often hides: the cost of a computation depends on
the grammar in which the computation is performed. A rule universe has
state counts, color counts, tape conventions, input encodings, halting
criteria, and output conventions. Change any of these and the observed
microcosm can change. The τ-framework makes the analogous dependency
explicit on its own side by refusing to speak about efficient search
until the agent, register grammar, interface width, and witness type have
been specified.

There is also an important sociological virtue here. The finite
microcosms are inspectable. A reader can see particular machines, runtime
tables, path graphs, and exception cases. That inspectability is why the
essay is a better anchor for this note than a purely philosophical essay
about computational irreducibility would be. It gives the comparison an
object: not merely ``hardness exists,'' but ``here is how hardness appears
when a finite rule class is made explicit.''

\subsection{Outliers, isolates, and the failure of smoothness}

One of the essay's strongest findings is qualitative rather than merely
enumerative. Runtime profiles in arbitrary small machines are not smooth
objects. They can jump. They can contain exceptional inputs. They can
depend on whether a machine halts, whether a larger machine is allowed, or
whether one changes the amount of nondeterminism. For constructed
algorithms, runtime often looks regular because the algorithm was designed
to have a regular structure. For machines in the wild, the behavior is
messier \cite{Wolfram2026RuliologicalPNP}.

The ``isolate machine'' idea sharpens this point. If a machine is the only
machine in a fixed finite class that computes a function, then its runtime
profile becomes a class-relative lower bound for that function within that
class. But even here the conclusion remains bounded. The same function may
be computed differently by a larger machine. The isolate shows that a
finite computational microcosm can contain genuine non-shortcut structure;
it does not license an unbounded asymptotic conclusion
\cite{Wolfram2026RuliologicalPNP}.

The isolate-machine surface is especially useful for a τ comparison
because it separates two forms of uniqueness. In Wolfram's setting,
uniqueness can mean uniqueness of a rule within a finite enumeration
class. In the τ setting, uniqueness usually means uniqueness of a
canonical address or normal representative after a construction grammar
has been applied. These are not the same kind of uniqueness. The former
is empirical and class-relative; the latter is structural and internal.
The bridge is that both resist the idea that all computations are
fungible black boxes. The boundary is that one uniqueness claim lives in
an enumerated rule census, while the other lives in an address-resolution
theorem.

Runtime outliers play a related role. A smooth complexity curve tempts
one to extrapolate. An outlier interrupts the extrapolation and asks
which hidden structure was missed. Wolfram's examples make that
interruption visible at small scale. The τ-framework then asks a sharper
question: is the outlier a sign of missing finite address structure, or
is it merely a bad representation of an admissible computation? That
question is not answered by the Wolfram essay, but the essay makes the
question vivid \cite{Wolfram2026RuliologicalPNP}.

\subsection{Nondeterminism as reach}

The P-vs-NP relevance becomes explicit when Wolfram introduces
nondeterministic and multiway machines. A deterministic machine applies one
rule at each step. A nondeterministic machine allows a choice among rules,
creating a multiway graph of possible paths. In the essay's finite
examples, this can shorten computations substantially. Some functions that
take exponential time for a small deterministic machine can be reached much
faster by a small nondeterministic machine. Other hard deterministic cases
do not succumb to the smallest nondeterministic augmentation. That
non-uniformity is one of the most useful facts for the present note
\cite{Wolfram2026RuliologicalPNP}.

Nondeterminism, in this setting, should be read as reach. It allows paths
through a multiway graph that deterministic rules do not follow. But reach
is not yet construction. A path may exist without giving a deterministic
procedure for finding it. A multiway graph may contain a short route
without making that route visible to a bounded deterministic observer. That
distinction is the hinge on which the comparison with τ-admissibility
turns \cite{Wolfram2026RuliologicalPNP}.

This is also where finite experiments become philosophically sharp.
Inside a small machine class, nondeterminism is not an abstraction
floating above the model. It is a concrete addition of rule alternatives.
One can draw the branch graph. One can ask which paths halt, which paths
compute the same output, and which paths shorten the route. That
concreteness prevents a common confusion: nondeterminism is not magic
knowledge. It is an expanded path structure. Whether that path structure
can be converted into deterministic construction is a further question.

The τ note takes that further question as central. It asks not only
whether there exists a short accepting branch, but whether the object
certified by that branch has finite constructive coordinates. If the
witness is still a black-box certificate, nondeterminism has supplied
reach but not structure. If the witness decomposes into per-prime
components and reassembles canonically, the story changes. That is why
the τ comparison cannot stop at the word ``speedup.'' It has to ask what
kind of thing has been sped up.

\subsection{Everything-machine and ruliad limit}

The essay then considers the limiting case of adding all conceivable rule
choices inside a finite class: the ``everything machine.'' This is a
maximally nondeterministic object. It can reach outputs in the minimum
number of steps structurally possible, and it functions as a fragment of
the ruliad. In Wolfram's broader vocabulary, the ruliad is the entangled
limit of all possible computations. It includes all possible computational
paths, while actual observers sample bounded slices of it
\cite{Wolfram2026RuliologicalPNP}.

This is the deepest philosophical bridge to the Panta Rhei note, and also
the easiest place to overclaim. The everything-machine/ruliad endpoint is
about maximal reach. The τ computation spine is about admissible
construction under a specific address grammar. Both make observer bounds
and rule-space structure visible. But their primitives differ. The ruliad
is a maximal computational ensemble; \CategoryTau\ is a constructed
framework with a primorial tower, CRT decomposition, finite interface
width, and a product-meet proof architecture. The bridge is real because
both ask what computation looks like when rule space is made explicit. The
bridge is limited because maximal reach is not the same as finite
admissible construction.

The everything-machine idea is nevertheless valuable because it provides
a limiting counter-image to the τ collapse. In one direction, one tries
to include more and more possible rules until reachability becomes
maximal. In the other, one tries to identify the finite quotient through
which a particular computation stabilizes. These are almost opposite
gestures. The ruliological gesture is expansive: include more rules. The
τ-admissible gesture is resolving: identify the finite structure that is
actually needed. Their proximity lies in the fact that both reject a
single fixed black-box machine as the final object of study. Their
difference lies in the direction of explanation.

\subsection{The essay's own boundary}

The essay's conclusion is unusually useful for our purposes because it
keeps the orthodox boundary intact. Wolfram states that finite
ruliological subquestions do not directly resolve the full P-vs-NP
problem. They give intuition and reveal difficulties: runtime wildness,
halting subtleties, limit-order issues, and the role of computational
irreducibility. That is exactly the boundary this note must preserve. The
Wolfram essay is an anchor because it makes the shape of the difficulty
visible; it is not evidence that any external theorem has been settled
\cite{Wolfram2026RuliologicalPNP}.

This boundary is one reason the essay can be used responsibly. If the
essay had claimed to solve P vs NP outright, the present note would first
have to adjudicate that claim. Because it does not, the note can use the
essay for what it does exceptionally well: it shows how quickly finite
rule spaces become rich enough to resist simple summarization. That is a
better bridge to τ-admissibility than a premature proof claim would be,
because τ-admissibility is itself a boundary concept. It says: here is
the fragment where construction is available; outside it, do not pretend
that search has already been tamed.

% =======================================================================
\section{The \texorpdfstring{$\tau$}{tau} computation layer}
\label{sec:tau-layer}
% =======================================================================

The τ computation spine begins with a typing claim. P-vs-NP-style
questions are not primitive at every layer of the framework. They require
agents, steps, and an efficiency measure. Without an agent there is no
searcher. Without steps there is no runtime. Without a relationship
between input size and step count there is no polynomial-versus-exponential
distinction. The computation layer says that these ingredients become
native only at \Etwo, where codes contain their own decoders and execution
is a Code--Execution--Code cycle \cite{PR-III}.

This is not an orthodox complexity theorem. It is a framework typing
statement. It says where the question becomes meaningful in the internal
architecture. At \Ezero, the framework has addresses and static structure,
but not processes. At \Eone, it has dynamics, but not self-referential
codes. At \Etwo, it has agents whose decoder is itself representable as a
code and whose execution produces new codes inside the same operational
layer. Only there can one ask whether a search procedure is efficient.

This typing move matters for the comparison with Wolfram. Wolfram begins
with finite rule systems and asks what their computed functions and
runtimes look like. The τ source spine begins one level earlier: what
structure must exist before a runtime question has a subject? Once the
subject exists, the framework does not ask only whether a nondeterministic
branch exists. It asks whether the witness lives inside a finite address
architecture that a deterministic τ-agent can construct.

The resulting picture is less dramatic than a headline, but more useful
as a research claim. The framework does not say that an arbitrary
mathematical question becomes computationally easy once translated into
τ-language. It says that the computational version of the question has a
type. Before \Etwo, there may be structures, transformations, and
dynamics, but not yet a self-referential computational agent with an
internal decoder. At \Etwo, a program can act on codes and return codes
inside the same typed layer. That is the first point at which the
phrases ``solver,'' ``verifier,'' and ``witness'' have their native
framework meaning.

This also prevents a false comparison with Wolfram's small machines.
Wolfram's machines are deliberately minimal. They show how much behavior
can be generated by tiny rule systems. The τ computation layer is not
minimal in that sense; it is already a typed operational layer built on
earlier address and enrichment structure. Its strength is not that it
uses fewer primitives. Its strength is that it can say which primitives
are doing the computational work. A finite Turing-machine enumeration
finds difficulty by varying rules. The τ spine locates difficulty by
asking whether the computation has earned finite readout.

\begin{widetable}[t]
\caption{The τ computation/admissibility source spine used in this note.
Book labels are source provenance; the manuscript frames the material as
the τ computation spine.}
\label{tab:tau-spine}
\prtablefont
\begin{tabularx}{\textwidth}{p{0.18\textwidth}p{0.17\textwidth}Y Y}
\toprule
\textbf{Source locus} & \textbf{Registry IDs} & \textbf{Claim used here} & \textbf{Guardrail} \\
\midrule
Computation layer & \regid{III.D49}, \regid{III.D50}, \regid{III.R22}
& Agents, steps, and efficiency become native at \Etwo. & Typing claim, not orthodox theorem. \\
τ-Tower Machine & \regid{III.D51}, \regid{III.T30}, \regid{III.D52}
& TTMs have finite control over τ-addressed registers and bounded observable transitions. & TTMs are not Wolfram machines. \\
Interface width & \regid{III.D53}, \regid{III.D54}, \regid{III.T31}
& τ-admissibility is finite stabilization in the primorial tower. & Applies inside \Adm. \\
Witness search & \regid{III.D55}, \regid{III.P22}, \regid{III.P23}
& Witnesses are τ-addresses decomposed by CRT into per-prime components. & Polynomial in primorial depth \(k\), not automatically bit length. \\
Computational bi-square & \regid{III.D56}, \regid{III.T32}, \regid{III.T33}
& Product-Meet Collapse yields \(\Padm=\NPadm\). & Not classical \(\ClassP=\ClassNP\). \\
No barrier / ZFC horizon & \regid{III.T34}, \regid{III.D79}, \regid{III.R46}
& The internal encoding gap is removed, while the ZFC bridge remains horizon-limited. & Provability/export language is scoped and partly conjectural. \\
\bottomrule
\end{tabularx}
\end{widetable}

% =======================================================================
\section{\texorpdfstring{$\tau$}{tau}-Tower Machines and interface width}
\label{sec:ttm-interface}
% =======================================================================

The τ-Tower Machine is the concrete machine grammar used by the internal
source spine. It resembles a register machine in having finite control and
finitely many registers. It differs because the registers hold
τ-addresses, not ordinary bitstrings. A program, its input data, its
instruction set, and its output are all internal address-bearing objects.
The observable transition of a machine reads a bounded window of this
address structure, which is why constant-width tableau language appears in
the source spine.

This is not a cosmetic change of alphabet. In an ordinary Turing model,
the tape is a flat string of symbols and the input-size metric is tied to
that representation. In a TTM, the relevant data are organized by the
primorial tower and by the finite quotient through which the computation
stabilizes. The key invariant is \emph{interface width}: the amount of
primorial depth a computation must inspect before deeper levels contribute
no new information.

\begin{definition}[Interface-width reading \scopetau]
For this note, a computation is treated as τ-admissible only when its
verifier has finite interface width: there is a depth \(k_0\) such that,
for all relevant input sizes, the verifier's observable behavior factors
through the finite quotient at depth \(k_0\). The infinite tower then
collapses, for that computation, to a finite computational horizon.
\end{definition}

The phrase ``for that computation'' is load-bearing. It blocks the most
tempting overclaim. The source spine does not say that every arbitrary
problem in every formal system magically has small interface width. It
defines a class of admissible computations and proves the internal
consequences inside that class. A function whose required depth grows with
the input remains outside the admissible fragment until a finite-width
factorization is supplied.

Interface width should therefore be read as a test, not a decoration. A
candidate τ computation must answer concrete questions. What finite
primorial quotient carries the relevant information? What does the
verifier inspect? Which deeper coordinates are provably irrelevant to the
accept/reject decision? How is the stability witnessed? If these
questions cannot be answered, the computation has not entered the
admissible fragment merely because it has been described in τ language.

The width condition also explains why the internal complexity parameter
is not automatically ordinary input length. Primorial depth \(k\)
controls a tower-indexed finite quotient. Ordinary bit length controls
the encoding size of a string in a Turing representation. Relating the
two requires an export theorem or at least an explicit encoding policy.
This note uses \(k\)-polynomial language only in the internal sense. It
does not silently convert it into a classical polynomial-time statement.

\begin{definition}[τ-admissible witness problem \scopetau]
For this note, a τ-admissible witness problem is a decision problem
presented to an \Etwo-level τ-agent such that:
\begin{enumerate}[leftmargin=1.2em,itemsep=1pt,topsep=2pt]
\item the verifier factors through a finite primorial quotient;
\item accepting witnesses are τ-addresses at that quotient;
\item the witness constraints decompose into finite local conditions on
the relevant prime-indexed coordinates;
\item compatible local witnesses reassemble canonically by CRT; and
\item the verifier's deeper tower data do not change the accept/reject
decision.
\end{enumerate}
If any of these clauses is missing, the problem has not yet entered the
admissible fragment used by the collapse claim.
\end{definition}

\begin{definition}[The classes \texorpdfstring{\(\Padm\)}{tau-Padm} and \texorpdfstring{\(\NPadm\)}{tau-NPadm} \scopetau]
\(\Padm\) denotes the τ-admissible decision problems for which a
deterministic τ-Tower Machine constructs the accepting address, when it
exists, in time polynomial in the internal primorial depth. \(\NPadm\)
denotes the τ-admissible decision problems for which a finite-width
τ-verifier checks an accepting address in polynomial internal depth. Both
classes are already restricted by τ-admissibility; neither is the
ordinary \(\ClassP\) or \(\ClassNP\) class of orthodox complexity theory.
\end{definition}

The comparison with Wolfram is therefore precise. Wolfram fixes finite
machine classes and studies the functions and runtimes that appear within
them. The τ spine fixes an internal computation grammar and asks whether a
verifier stabilizes through a finite quotient. Both approaches turn an
infinite question into finite surfaces one can inspect. But the surfaces
are different: finite enumeration in one case, finite interface width in
the other.

This difference suggests a productive diagnostic. If a Wolfram-style
finite example looks hard because no smaller machine computes it in the
chosen class, one should not immediately ask whether it is τ-admissible.
One should first ask a weaker question: what is the analogue of its
interface? What information must any verifier read to certify the output?
Does the difficulty come from a large search space, from an unstable
representation, from halting sensitivity, or from the absence of a finite
factorization? These questions are not solved in this note, but they
turn the bridge into a usable research protocol.

% =======================================================================
\section{A one-minute toy example: assembled witness, not scan}
\label{sec:toy-example}
% =======================================================================

The smallest useful example is deliberately not an NP-complete problem,
not a SAT reduction, and not an orthodox complexity reduction. It only
illustrates the τ mechanism: when local finite coordinates and a
reconstruction theorem are present, witness existence can become witness
assembly. Work in the finite quotient
\[
  Q=\mathbb{Z}/30\mathbb{Z}
  \cong
  \mathbb{Z}/2\mathbb{Z}\times
  \mathbb{Z}/3\mathbb{Z}\times
  \mathbb{Z}/5\mathbb{Z}.
\]
Suppose the local witness constraints are
\[
  \begin{aligned}
  x &\equiv 1 \pmod 2,\\
  x &\equiv 2 \pmod 3,\\
  x &\equiv 4 \pmod 5.
  \end{aligned}
\]
A flat scan would test residues \(0,\ldots,29\) until one satisfies all
three checks. CRT instead constructs the address. The three standard
reassembly terms are \(15\), \(10\), and \(6\): each is \(1\) on its own
prime component and \(0\) on the other two after the appropriate local
normalization. Thus
\[
  x \equiv 1\cdot 15 + 2\cdot 10 + 4\cdot 6
  =59 \equiv 29 \pmod {30}.
\]
The result can be checked immediately:
\[
  \begin{aligned}
  29&\equiv 1 \pmod 2,\\
  29&\equiv 2 \pmod 3,\\
  29&\equiv 4 \pmod 5.
  \end{aligned}
\]

This toy example captures only the address-construction grammar. It does
not prove anything about arbitrary SAT instances, ordinary NP-complete
problems, or the Clay P-vs-NP problem. Its purpose is narrower and more
useful: it shows the difference between scanning a flat finite space and
constructing a global witness from typed local components. In the
τ-admissible case, the source-spine theorem claims that the real witness
spaces have an analogous finite quotient, local constraint, and canonical
reassembly structure. When that structure is absent, the example does not
apply.
The work was moved from scanning candidate witnesses to constructing the
coordinate system in which the witness is already locally determined.

The failure variant is just as important. If the local constraints are not
independent, if the moduli are not compatible, or if a hidden global
consistency condition remains after local checking, the CRT assembly
pattern fails. Such a problem may remain a search problem and is not
τ-admissible by this example alone.

% =======================================================================
\section{Witness search as address resolution}
\label{sec:witness-search}
% =======================================================================

In the orthodox language of NP, a witness is a certificate. It may be easy
to check once supplied, even if it is hard to find. The usual anxiety is
that the witness space is too large: exponentially many candidate
bitstrings, no general deterministic method for locating the right one,
and no known proof that such a method cannot exist.

The τ source spine retypes the witness. A witness for an admissible
problem is a τ-address at some primorial depth. That address carries
coordinates. Its components can be read through the CRT decomposition of
\(\mathbb{Z}/\Prim(k)\mathbb{Z}\) into per-prime residues. Under
τ-admissibility, the verifier factors through finite width, so witness
constraints become per-prime constraints that can be searched locally and
reassembled globally. The witness is no longer a flat object hidden in a
featureless space. It is a structured address whose local components have
to cohere.

This is the technical heart of the title. Search becomes structure not
because nondeterminism is wished away, but because the object searched for
has changed type. In a flat bitstring space, global satisfaction is a meet
over constraints with no guaranteed product architecture. In the primorial
tower, when the admissibility hypotheses hold, the CRT turns global
satisfaction into compatible local construction. The per-prime witness
spaces remain finite, and their reassembly is unique.

It is useful to spell this out as a four-step grammar. First, the
verifier is typed: it is not an arbitrary external checker but an
\Etwo-level computation whose observable behavior has a finite interface.
Second, the candidate witness is typed: it is represented as an address
at a finite primorial depth rather than as an undifferentiated string.
Third, the constraints are localized: acceptance conditions factor into
conditions on finite per-prime components. Fourth, the local solutions are
reassembled by CRT into the global address. At no point is the witness
found by a mysterious act of guessing. The work is shifted from blind
global enumeration to structured local resolution.

This is also why the address language is not merely metaphorical. In an
ordinary NP problem, a certificate can be encoded as bits, but the bit
encoding does not by itself explain why the right certificate should be
found efficiently. It only gives a representation in which verification
is possible. In the τ-admissible case, the representation is tied to a
construction theorem: the finite quotient, per-prime decomposition, and
unique reassembly are part of the computational grammar. The address is
not a label attached after the fact. It is the shape through which the
witness becomes constructible.

\begin{proposition}[Bridge reading of witness search \scopetau]
Within the τ-admissible fragment, witness search is not modeled as an
unstructured enumeration of certificates. It is modeled as address
resolution: identify the finite quotient, solve the per-prime witness
conditions, and reassemble the canonical address by CRT.
\end{proposition}

\begin{proof}[Source-spine proof sketch]
The source spine supplies three ingredients. First, finite interface width
forces the verifier to factor through a finite primorial quotient. Second,
the witness space at that quotient decomposes into per-prime components by
CRT. Third, unique CRT reassembly turns compatible local witnesses into a
global witness. The cost statement in the source is polynomial in the
primorial depth \(k\), with the important caveat that this is not
automatically the same as ordinary bit-length polynomiality.
\end{proof}

The contrast with Wolfram's nondeterministic machines is now clear.
Wolfram's multiway graphs show that branch availability can shorten paths.
The τ source spine asks for something stronger: a deterministic structural
procedure that finds and reassembles the witness because the witness space
has finite address architecture. A short branch and a constructed address
are related ideas, but they are not the same.

\subsection{What the toy example did and did not show}

The CRT example above is intentionally small enough to be checked by
hand. It shows the contrast between a flat search space and an address
space whose local components can be assembled. It does not show that
every certificate space has such a product structure. It also does not
show that ordinary input length and primorial depth are interchangeable.
Those are precisely the questions that any classical export would have
to answer.

The didactic point is therefore conditional. If a witness problem has a
finite quotient, local constraints, and canonical reassembly, then the
right search grammar is construction. If the witness remains a flat
certificate with no finite address decomposition, the example has no
force. This is why the definitions of τ-admissibility come before the
formal equality in this note.

% =======================================================================
\section{The computational bi-square and admissibility collapse}
\label{sec:bisquare-collapse}
% =======================================================================

The computational bi-square is the diagrammatic center of the internal
claim. Earlier bi-squares in the framework have the same repeated shape:
one square records tower coherence, another records spectral naturality,
and the pasted rectangle expresses the relevant constraint at that layer.
At the computation layer the objects have changed. They are verifier
computations and witness spaces. The left square is the execution square;
the right square is the witness square. Pasting them yields the
Product-Meet Collapse.

In plain language, the Product-Meet Collapse says that, inside the
admissible setting, satisfying all witness constraints is realized by
constructing the product of compatible per-prime witness components. The
meet of constraints becomes the product of components. This is the moment
where search becomes structure in the strongest internal sense.

The diagram matters because it prevents the collapse from becoming a
black-box slogan. The left side records that computation is still an
execution process: an agent acts, a verifier reads a finite interface,
and the observable transition must remain coherent through the tower. The
right side records that the witness is not merely asserted: its local
components have to be natural with respect to the spectral/prime-indexed
decomposition. The pasted diagram says that these two requirements are
not independent decorations. They meet in the same finite construction
surface.

This is also where the τ source spine differs from a lookup-table
shortcut. In a finite universe, one can sometimes compute a function by
encoding enough answers in a larger machine. That is a representation
trick unless it comes with a stable construction rule. Product-Meet
Collapse is not a claim that every finite instance can be stored. It is a
claim that, for admissible computations, the witness space itself has a
decomposition that the solver can use uniformly. The distinction is
essential for any future export discussion.

\begin{widefigure}[t]
\begin{tikzpicture}[
  node distance=0.8cm and 0.55cm,
  box/.style={draw=pr-brand-700, rounded corners=2pt, thick,
    align=center, inner sep=5pt, text width=0.18\textwidth,
    fill=pr-brand-50},
  taubox/.style={draw=pr-ink-700, rounded corners=2pt, thick,
    align=center, inner sep=5pt, text width=0.18\textwidth,
    fill=white},
  labelbox/.style={font=\sffamily\bfseries\small, text=pr-brand-700},
  arrow/.style={-{Latex[length=2.2mm]}, thick, color=pr-ink-700}
]
\node[labelbox] (lefttitle) {Finite ruliological exploration};
\node[box, below=0.45cm of lefttitle] (rules) {fixed finite\\rule classes};
\node[box, right=of rules] (profiles) {computed functions\\and runtime profiles};
\node[box, right=of profiles] (hardness) {outliers, isolates,\\multiway speedups};
\node[box, right=of hardness] (pressure) {irreducibility\\pressure};

\node[labelbox, below=2.1cm of lefttitle] (righttitle) {τ-admissible construction};
\node[taubox, below=0.45cm of righttitle] (agent) {\Etwo\ agent\\and τ-addresses};
\node[taubox, right=of agent] (width) {finite interface\\width};
\node[taubox, right=of width] (witness) {CRT witness\\decomposition};
\node[taubox, right=of witness] (collapse) {Product-Meet\\Collapse};

\draw[arrow] (rules) -- (profiles);
\draw[arrow] (profiles) -- (hardness);
\draw[arrow] (hardness) -- (pressure);
\draw[arrow] (agent) -- (width);
\draw[arrow] (width) -- (witness);
\draw[arrow] (witness) -- (collapse);

\draw[arrow, dashed] (hardness.south) -- node[right, font=\scriptsize\itshape,
  text=pr-ink-600, align=center] {comparison,\\not identity} (witness.north);
\draw[arrow, dashed] (pressure.south) -- node[right, font=\scriptsize\itshape,
  text=pr-ink-600, align=center] {obstruction\\vs criterion} (collapse.north);
\end{tikzpicture}
\caption{Search becomes structure only on the lower path. Wolfram's finite
microcosms make hardness surfaces visible from below. The τ path turns
search into construction only after finite interface width and CRT witness
decomposition are available.}
\label{fig:search-becomes-structure}
\end{widefigure}

\begin{theorem}[Framework source-spine theorem: τ-admissibility collapse \scopetau]
For τ-admissible witness problems, the internal source spine states
\[
  \Padm = \NPadm .
\]
If a verifier has finite interface width, witness search runs in time
polynomial in the primorial depth used by the admissible computation.
\end{theorem}

\begin{proof}[Source-spine proof sketch]
The source-spine proof assembles three claims. Constant-width verification
supplies the execution side. CRT witness decomposition supplies the
witness side. Product-Meet Collapse pastes the two squares: independently
resolved per-prime witnesses assemble into a global witness. Thus
verification and search both become polynomial in the internal depth
parameter for the admissible fragment.
\end{proof}

The theorem remains source-spine material \cite{PR-III}. The surrounding
public construction grammar can be inspected through the Construction
Spine, TauLib, WP003, the Release Manifest, and the Corpus Registry
\cite{PR-ConstructionSpinePublic,PR-TauLibCorpusPublic,PR-TauLibWhitePaperPublic,PR-ReleaseManifestPublic,PR-RegistryPublic}.

\prcallout[Main non-claim]{%
The displayed equality is \(\Padm=\NPadm\), not an orthodox
\(\ClassP=\ClassNP\) theorem. It belongs to τ-admissible computation.
This is an internal equality for τ-admissible witness problems. It does
not settle orthodox P vs NP over unrestricted Turing-machine / ZFC
encodings.}

% =======================================================================
\section{What fails \texorpdfstring{$\tau$}{tau}-admissibility?}
\label{sec:admissibility-failures}
% =======================================================================

The collapse claim is credible only if the failure modes are visible. A
problem does not become τ-admissible because it has been renamed in
τ-vocabulary. It has to earn the finite construction surface. The most
important failure modes are:

\begin{widetable}[t]
\caption{Failure modes for τ-admissibility. The restrictive nature of
this table is what prevents \(\Padm=\NPadm\) from being an unguarded
\(\ClassP=\ClassNP\) claim.}
\label{tab:admissibility-failures}
\prtablefont
\begin{tabularx}{\textwidth}{p{0.23\textwidth}Y Y}
\toprule
\textbf{Failure mode} & \textbf{Why τ-admissibility fails} & \textbf{Consequence} \\
\midrule
Unbounded witness representation & The witness cannot be finitely
address-coded at the cutoff. & The problem remains search. \\
Hidden global consistency & Local checks do not reconstruct the global
witness. & Additional global search remains. \\
No reconstruction theorem & No CRT-, Yoneda-, or Hartogs-style assembly
route is available. & Existence has not become construction. \\
Unbounded interface width & Local data require growing cross-talk as
instances vary. & The verifier is not boundedly admissible. \\
External semantic oracle & Verification depends on meaning outside formal
address data. & Runtime or semantics have been hidden. \\
Encoding-dependent shortcut & The collapse depends on changing the
problem representation illegitimately. & The result is not exportable. \\
\bottomrule
\end{tabularx}
\end{widetable}

These are not embarrassment clauses; they are part of the theorem's
meaning. A serious future benchmark should contain negative rows where
the admissibility test fails. Wolfram-style finite microcosms are useful
here because they can produce exactly the kind of stubborn runtime
outliers and branch structures that test whether a proposed τ encoding is
real structure or only a renamed search.
This is an internal equality for τ-admissible witness problems. It does
not settle orthodox P vs NP over unrestricted Turing-machine / ZFC
encodings.

% =======================================================================
\section{Finite exploration vs finite interface width}
\label{sec:finite-interface}
% =======================================================================

We can now state the first major comparison. Wolfram uses finite
enumeration. The τ spine uses finite interface width. Both are ways of
refusing to let the infinite formulation do all the work.

Finite enumeration starts by fixing a rule class and asking what appears
inside it. One can count machines, compare functions, identify isolated
representatives, and inspect runtime profiles. The gain is empirical
specificity. The cost is class relativity. A fact about \(s=3,k=2\)
machines may change when \(s=4,k=2\) machines are admitted.

Finite interface width starts differently. It asks whether a verifier
stabilizes through a finite quotient independent of deeper tower levels.
The gain is theorem-grade internal structure. The cost is admissibility:
only computations satisfying the finite-width criterion fall inside the
collapse. A computation that truly requires unbounded depth remains outside
until additional structure is supplied.

\begin{widetable}[t]
\caption{Core comparison axes. The table is intentionally asymmetric:
Wolfram supplies finite microcosms of difficulty; the τ spine supplies a
conditional construction theorem under admissibility.}
\label{tab:comparison}
\prtablefont
\begin{tabularx}{\textwidth}{p{0.18\textwidth}Y Y Y}
\toprule
\textbf{Axis} & \textbf{Wolfram ruliology} & \textbf{τ computation spine} & \textbf{Bridge sentence} \\
\midrule
Finite structure & Enumerated finite machine universes. & Finite interface width and finite quotient factorization. & Both make finite readout explicit. \\
Search & Runtime outliers and class-relative difficulty. & Witness search becomes address resolution under τ-admissibility. & Wolfram studies hard search; τ specifies when search becomes construction. \\
Nondeterminism & Multiway branching can shorten paths. & Witnesses are canonical addresses. & Reach and construction must be kept distinct. \\
Irreducibility & Computational irreducibility obstructs shortcutting. & Inadmissibility marks failure of finite collapse. & Irreducibility is pressure; admissibility is criterion. \\
Infinity & Ruliad/everything-machine gestures toward maximal rule-space reach. & τ-native computation uses a particular primorial construction grammar. & Philosophically adjacent, architecturally different. \\
Orthodox export & Essay does not settle classical P vs NP. & τ collapse is internal unless a separate bridge is supplied. & Shared humility makes the bridge note viable. \\
\bottomrule
\end{tabularx}
\end{widetable}

This is why the note's title should not be read as a universal claim that
every search becomes structure. The statement is conditional. Search
becomes structure when a finite address architecture exists, when the
verifier stabilizes at finite interface width, and when the witness space
factorizes in a way compatible with construction. Wolfram's finite
microcosms help us see why those conditions matter: without them, runtime
behavior can remain wild, class-relative, and irreducibly difficult to
summarize.

\subsection{Two finite moves, two different risks}

The two finite moves have opposite failure modes. Finite enumeration can
be too weak because it remains tied to the chosen class. A finite machine
microcosm may contain a genuine outlier, but another class may remove it.
The result is empirical richness without immediate generality. Finite
interface width can be too strong because it excludes computations whose
observable behavior does not stabilize. The result is theorem-grade
structure, but only inside a typed fragment.

The future research program has to keep both risks visible. If the τ
side ignores Wolfram-style finite exploration, it risks becoming too
formal: internally elegant but insufficiently tested against the messy
behavior of small rule universes. If the Wolfram side ignores finite
interface conditions, it risks staying descriptive: rich catalogs of
difficulty without a criterion for when a witness space has become
constructible. The bridge note is valuable precisely because it can let
each side expose the other's blind spot.

This is also why the comparison should be carried out with examples,
not only with vocabulary. A finite rule example can be annotated by the
questions it raises for τ: what is the interface, what is the witness
object, what would count as a finite quotient, and what would fail if no
such quotient exists? Conversely, a toy τ-admissible example can be
translated into a Wolfram-style finite exploration: what does the
deterministic search look like before the address decomposition is made
visible? The two directions would not prove equivalence. They would make
the bridge measurable.

% =======================================================================
\section{Nondeterminism versus witness construction}
\label{sec:nondeterminism}
% =======================================================================

The second comparison is the most delicate. In Wolfram's essay,
nondeterminism is implemented by allowing a choice among rules at each
step. The resulting multiway graph contains many possible paths. Some
paths compute outputs faster than the deterministic rule path. In the
limiting everything-machine picture, nondeterminism becomes maximal reach:
from an input, many outputs become accessible by structurally minimal
paths.

In orthodox NP language, nondeterminism is often described by saying that
the right witness can be guessed. That phrase is useful pedagogically, but
it hides the key issue. A short accepting path may exist without a
deterministic polynomial method for finding it. A multiway graph may
contain the route, but the route may not be visible to a bounded solver.
The real question is not merely whether the path exists. It is what kind
of structure makes the path constructible.

The τ answer is not ``branch harder.'' It is ``change the type of the
witness.'' Under admissibility, the witness has address structure. The
search space is not just a cloud of possibilities; it is a finite
quotient with per-prime components. The deterministic construction
procedure is not a blind traversal of branches. It is a decomposition and
reassembly procedure.

\begin{remark}[Reach is weaker than construction]
A nondeterministic reachability statement says that some path exists in a
multiway graph. A τ-admissible construction statement says that the
witness can be recovered from finite address components and reassembled.
The second is stronger and more typed than the first.
\end{remark}

This distinction also protects the note from an attractive but false
symmetry. It would be easy to say that Wolfram's nondeterministic speedups
are the same thing as τ witness construction. They are not. Wolfram's
examples show how branch availability can reduce observed path length in
finite rule universes. The τ spine gives an internal criterion for when a
witness space ceases to be featureless and becomes constructible. The
bridge is that both focus on the structure of possible paths. The
boundary is that only the τ side claims an internal construction theorem,
and only under admissibility.

The distinction has a practical consequence for future benchmarks. A
candidate example should not be counted as a τ analogue merely because
nondeterminism speeds it up. The first diagnostic should be: does the
speedup identify a reusable witness structure, or only a fortunate
branch? If the shortest branch depends on global search through the
multiway graph, then the example remains a reachability example. If the
branch can be replaced by a local decomposition-and-reassembly procedure,
then it starts to resemble the τ witness grammar.

This is a useful discipline because many popular accounts of P vs NP
slide too quickly between ``the answer can be checked'' and ``the answer
can be found.'' Wolfram's finite graphs make the slide visible: the path
can be in the graph without being efficiently discoverable. The τ source
spine turns that visibility into a criterion: discoverability comes from
finite address structure, not from the mere existence of a branch.

% =======================================================================
\section{Irreducibility, inadmissibility, and infinity}
\label{sec:irreducibility}
% =======================================================================

Wolfram's essay repeatedly returns to computational irreducibility: the
possibility that no shortcut captures the behavior of a computation
without effectively running it. In the P-vs-NP context, this pressure
appears as runtime wildness, halting uncertainty, difficult limit behavior,
and the suspicion that exceptions may not average away. The essay suggests
that computational irreducibility and undecidability pressure are not rare
corner cases but pervasive features of rule space.

The τ source spine uses a different vocabulary. It does not begin by
asserting irreducibility. It defines admissibility. A computation is
admissible when finite interface width exists. A computation is
inadmissible, for the purposes of this note, when no such finite
stabilization is available. The relationship is suggestive but not
identity. Irreducibility is a pressure against shortcutting. Inadmissibility
is a failure to meet a particular internal finite-width criterion.

This contrast helps both sides. Wolfram helps us resist a too-easy
reading of the τ collapse: if finite-width structure is absent, wildness
and non-shortcut behavior should be expected, not treated as a minor
technicality. The τ spine helps us sharpen the positive side: when
admissibility is present, the reason search becomes constructive is not
that all computations are reducible in general, but that this computation
has finite quotient structure.

The infinity comparison is similar. Wolfram's everything-machine and
ruliad language points toward maximal computational reach: all rules, all
paths, all computations, with bounded observers sampling slices. The τ
spine points toward a different treatment: a primorial tower whose finite
stages, potential limits, and address structure delimit what a τ-agent can
inspect and construct. Both frameworks make observer boundedness visible.
But maximal reach and admissible construction are not interchangeable.

This difference gives the note its negative testing surface. If a
candidate problem continues to require ever deeper primorial coordinates,
or if the verifier cannot be shown to factor through a stable finite
quotient, then the τ collapse should not be invoked. Such a case may be
computationally irreducible in a Wolframian sense, inadmissible in the τ
sense, both, or neither. The categories must not be collapsed without a
bridge argument.

The most interesting future cases may be mixed. A finite rule universe
may show irreducible-looking behavior that later turns out to have a
hidden address decomposition under a different representation. Conversely,
a τ-like encoding may appear promising until one discovers that the
required interface width grows with the input. The honest research
program is to sort these cases, not to declare them solved in advance.
That is why inadmissibility is not a failure of the framework; it is one
of its necessary boundary markers.

\subsection{Secondary echoes}

Book II supplies a finite/readout substrate echo: finite-stage data,
ultrametric depth, profinite limits, and stagewise finite linear algebra
\cite{PR-II}. Book VI supplies application echoes: evolution and cognition
as processes that navigate large search spaces through local exploration,
selection, caching, and approximation \cite{PR-VI}. Book VII supplies the
readout echo: apparent randomness and probability arise when observers
cannot resolve all internal structure \cite{PR-VII}. These echoes are
useful for later prose, but they are not the proof spine of this note.
They explain why the search/structure pattern recurs; they do not prove a
complexity-theoretic claim.

These secondary echoes are contextual only; they are not part of the
proof spine of this note. They can show that the same readout pattern
recurs across the framework: finite carriers before limit claims, local
navigation before global comprehension, observer resolution before
probabilistic summary. But they should not be allowed to blur the proof
spine. The P-vs-NP comparison lives or dies on the
computation/admissibility material, not on a general family resemblance
across the books.

% =======================================================================
\section{The ruliad is not \texorpdfstring{Category $\tau$}{Category tau}}
\label{sec:ruliad}
% =======================================================================

Because the two projects are philosophically close, the distinction must
be made explicitly. Wolfram's ruliad is a maximal object: the entangled
limit of all possible computations. It is not a single machine, not a
single model, and not a particular address grammar. It is the all-paths
object from which bounded observers sample slices.

\CategoryTau, by contrast, is a specific construction framework. It has
generators, axioms, a primorial tower, CRT decomposition, enrichment
levels, τ-addresses, and proof-organizing diagrams such as the
computational bi-square. Its internal P-vs-NP claim is not that all paths
exist. It is that, under admissibility, the witness space has enough
finite address structure for search to become construction.

The productive comparison is therefore this:

\begin{quote}\small
The ruliad foregrounds maximal computational possibility and observer
sampling. The τ computation spine foregrounds admissible construction and
finite readout. Both make the hidden architecture of computation visible,
but they do so with different primitives and different proof obligations.
\end{quote}

This distinction should stay visible in any public version of the note.
The temptation will be to say that the ruliad and \CategoryTau\ are two
names for the same intuition. They are not. They are neighboring
frameworks that converge on a shared dissatisfaction with black-box
complexity abstractions. That convergence is enough. It is more honest,
and more interesting, than forced equivalence.

One way to state the distinction is by asking what each framework treats
as primitive. The ruliad begins from the space of possible computations
and asks how observers are embedded in that enormous object. Category τ
begins from a constructed address calculus and asks how layers of
structure, readout, computation, and observer access are built. The
ruliad is maximal before observation; Category τ is constructive before
maximality. This is not a ranking. It is a difference in explanatory
direction.

Equivalently: the ruliad is a possible-computation space, while
\CategoryTau\ is a constrained constructive architecture. That difference
is exactly why the comparison is worth making. A possible-computation
space asks what exists in the maximal rule ensemble. A constructive
architecture asks what can be built, read, checked, and reassembled under
the finite constraints of the framework.

That difference matters for P vs NP. A maximal computational ensemble can
make the landscape of possible paths vivid, but it does not by itself
select a deterministic construction procedure. A constructed address
grammar can select a construction procedure, but only after its
admissibility hypotheses have been checked. The two perspectives are
therefore complementary in the ordinary sense: each can expose questions
the other tends to leave implicit. They are not complementary in the
stronger sense of forming two halves of a single proof.

% =======================================================================
\section{Public source and construction context}
\label{sec:public-context}
% =======================================================================

The τ side is grounded in the computation/admissibility source spine
cited as \cite{PR-III}, and the public Panta Rhei site exposes the
surrounding construction grammar through inspectable routes: the
Construction Spine, TauLib Corpus, TauLib verification route, WP003 TauLib
Technical Overview, Release Manifest, and Corpus Registry
\cite{PR-ConstructionSpinePublic,PR-TauLibCorpusPublic,PR-TauLibVerifyPublic,PR-TauLibWhitePaperPublic,PR-ReleaseManifestPublic,PR-RegistryPublic}.
The internal source spine is summarized publicly through the Construction
Spine, TauLib, WP003, the Release Manifest, and the Corpus Registry.
Those public routes are not substitutes for the technical source spine,
but they give the reader inspectable context for what the framework means
by construction, registry, formal route, and release discipline.

\paragraph{Public inspection surfaces}
The Construction Spine gives a public route into the staged build order.
The TauLib Corpus and Verify/TauLib pages expose the formalization and
inspection surface. WP003 records the TauLib technical overview and trust
budget. The Release Manifest pins the public verification state. The
Corpus Registry gives the object-indexed route into definitions,
theorems, structures, and dependencies. These surfaces are contextual
inspection routes, not substitutes for the internal source-spine theorem.

The Wolfram-side context is similarly layered. The primary external
anchor for this note is the January 2026 essay itself
\cite{Wolfram2026RuliologicalPNP}. The Wolfram Institute ruliology route
is used only as contextual support for the word ``ruliology'' and the
broader programmatic setting \cite{WolframInstituteRuliology}. It is not
used as a second proof source for the P-vs-NP argument.

This public-context split matters for outreach. The note can be sent to
a Wolfram/Wolfram Institute reader without asking them to accept private
terminology on trust: the public pages show the construction spine and
verification surfaces; the manuscript explains the guarded comparison;
and the claim boundary remains explicit at each step.

% =======================================================================
\section{Orthodox P-vs-NP guardrails}
\label{sec:guardrails}
% =======================================================================

The most important section of this note may be the one that says what it
does not do. The classical P-vs-NP problem is a problem in orthodox
complexity theory, formulated over Turing machines, polynomial time,
decision languages, encodings, reductions, and proof systems such as
\ZFC\ \cite{ClayPvsNP}. A framework-internal theorem about τ-admissible computation does not
automatically export to that setting. A finite ruliological experiment
does not settle an asymptotic theorem. A philosophical comparison between
rule-space exploration and address construction does not replace a proof.

The τ source spine is itself explicit about this boundary. It distinguishes
the internal admissibility collapse from the classical Turing-model
question. It treats the bridge to \ZFC\ as a separate issue, and parts of
the provability/independence discussion are horizon-limited or conjectural.
This note should preserve that humility because it is not a weakness. It
is the condition under which the comparison becomes credible.

\begin{widetable}[t]
\caption{Guardrail ledger. The same phrase can be safe or unsafe depending
on the tier in which it is used.}
\label{tab:guardrails}
\prtablefont
\begin{tabularx}{\textwidth}{p{0.18\textwidth}Y Y Y}
\toprule
\textbf{Tier} & \textbf{Safe statement} & \textbf{Risk} & \textbf{Required discipline} \\
\midrule
τ-internal & \(\Padm=\NPadm\) for τ-admissible computation. & Medium:
readers may drop the subscript. & Keep ``τ-admissible'' and ``internal''
near every strong claim. \\
τ-native / physical & Physical or τ-native extensions are framework
claims under stated admissibility assumptions. & High: sounds like a
headline about all computation. & Use only as later extension, not as
the note's main public claim. \\
ZFC / Turing / Clay & The orthodox bridge is not established by this
note. & Critical: readers infer a Millennium-problem claim. & State the
non-claim in introduction, theorem discussion, and conclusion. \\
Wolfram anchor & The essay is a finite ruliological comparator. & Medium:
anchor may be misread as validation. & Call it an anchor or comparator,
not evidence for τ. \\
\bottomrule
\end{tabularx}
\end{widetable}

There is a second guardrail, subtler but just as important. The note should
not hide behind internal vocabulary. If it uses \(\Padm=\NPadm\), it must
say in ordinary prose what the subscript does. It marks the admissible
fragment. It marks the finite-interface-width condition. It marks a
computational grammar different from the orthodox Turing tape. The
subscript is not decoration. It is the boundary.

\subsection{The three claim tiers}

The note keeps three tiers separate. The first tier is
\emph{τ-internal}: within the typed admissible fragment, the source spine
states \(\Padm=\NPadm\). The second tier is
\emph{τ-native/physical}: the framework may later interpret computation
inside physical or observer-relative readout settings, but those claims
inherit additional assumptions. The third tier is \emph{orthodox export}:
a statement about standard Turing-machine languages, polynomial time in
ordinary input length, reductions, and proof in a background system such
as \ZFC.

The present note uses the first tier, mentions the second only as future
context, and refuses the third as a current claim. This is not rhetorical
timidity. It is source fidelity. The internal theorem and the external
Clay problem are not the same proposition with different fonts. They have
different machines, different encodings, different size measures, and
different proof obligations.

The safest public wording is therefore not ``τ solves P vs NP.'' It is:
inside τ-admissible computation, witness search is retyped as address
construction, yielding \(\Padm=\NPadm\); the export to the orthodox
\(\ClassP\) versus \(\ClassNP\) problem remains a separate bridge
problem. This is an internal equality for τ-admissible witness problems.
It does not settle orthodox P vs NP over unrestricted Turing-machine /
ZFC encodings. That sentence is less dramatic. It is also the sentence
that can survive scrutiny.

% =======================================================================
\section{Future bridge tests}
\label{sec:future-tests}
% =======================================================================

The best ending for this bridge note is not triumph but work. The bridge
between finite ruliology and τ-admissible construction suggests concrete
next tests.

\paragraph{Finite benchmark ledger}
The first test is documentary: reproduce a small Wolfram-style finite
machine ledger and annotate it with the τ vocabulary. Which examples look
like pure runtime outliers? Which look like branch-reach examples? Which
have any plausible finite-interface-width analogue? This would not prove
anything about orthodox complexity, but it would make the comparison less
literary and more inspectable.

The ledger should record at least five fields: machine class, computed
function or finite input-output map, runtime profile, nondeterministic
shortening behavior if present, and the proposed τ diagnostic. The τ
diagnostic should be allowed to say ``none.'' That negative outcome would
be scientifically useful, because it would prevent the comparison from
becoming an all-purpose translation machine.

\paragraph{Witness-decomposition examples}
The second test is constructive: build toy witness spaces at small
primorial depth and show explicitly how CRT decomposition differs from
branch enumeration. A good example would display a flat candidate search,
the per-prime decomposition, and the unique reassembly step side by side.
The point would be to let readers see why ``reach'' and ``construction''
are different claims.

Subsequent benchmark examples should include the actual residues and the
actual reassembly calculation. The purpose is not to impress with
formalism. It is to let a skeptical reader check the claim by hand. A good
example should fit on one page: problem statement, flat witness space,
local residue constraints, CRT reassembly, and the explicit note that the
example is τ-internal rather than a classical reduction.

\paragraph{Irreducibility boundary map}
The third test is negative: identify examples where the finite-width
criterion fails or remains unknown. This is essential. A bridge note that
only shows positive collapse risks sounding like a universal shortcut
story. A serious research note should say where search does not yet become
structure.

The boundary map should distinguish at least three cases: known
admissible, known inadmissible, and unknown. It should also distinguish
``not yet encoded'' from ``encoded but width appears to grow.'' Those are
different failures. The first is a missing translation; the second is a
structural obstruction. Keeping them separate would make the note a
better research tool.

\paragraph{Orthodox export checklist}
The final test is a checklist for any future classical export. Such an
export would need to specify the encoding, the input-size metric, the
relation between primorial depth and bit length, the verifier model, the
reduction notion, and the proof system in which the claim is made. Until
that checklist is satisfied, the classical Clay problem remains outside
the claim.

This checklist is deliberately demanding. It includes: an encoding
from ordinary decision-language instances into τ-addressed objects; a
size comparison theorem relating bit length and primorial depth; a proof
that classical polynomial-time verifiers translate into finite-width
τ-verifiers, or a clear statement of the restricted class that does; a
deterministic construction algorithm whose cost is polynomial in the
exported size measure; and a proof environment in which the translation
is formalized. If any item is absent, the correct public statement is
``internal τ result with open export,'' not ``classical P equals NP.''

\begin{center}
\begin{minipage}{\columnwidth}
\captionof{table}{Export checklist for any future orthodox P-vs-NP
bridge. The present note does not satisfy this checklist; it records what
would have to be shown.}
\label{tab:export-checklist}
\prtablefont
\begin{tabularx}{\columnwidth}{p{0.38\columnwidth}Y}
\toprule
\textbf{Export requirement} & \textbf{Question} \\
\midrule
Encoding & Can orthodox instances be mapped to τ-address objects without
exponential blowup? \\
Verification preservation & Do τ-local checks correspond to ordinary
polynomial-time verification? \\
Reconstruction cost & Is reassembly polynomial in orthodox input length,
not only in primorial depth? \\
Uniformity & Is the construction uniform, rather than instance-specific
advice? \\
Semantic closure & Are all meanings carried by formal address data, with
no hidden semantic oracle? \\
Witness return & Does the constructed τ-witness translate back to a valid
orthodox witness? \\
\bottomrule
\end{tabularx}
\end{minipage}
\end{center}

\paragraph{Ruliological counter-tests}
A useful bridge should also invite counter-tests from the Wolfram side.
If finite machine exploration produces examples where nondeterministic
speedup appears robust across many expanded classes, but no finite
address decomposition can be found, those examples become pressure tests
for the τ admissibility criterion. Conversely, if a τ toy example has a
clean address decomposition, one should ask whether finite rule
enumeration can rediscover that structure or merely observe the speedup.
Both outcomes would teach something.

% =======================================================================
\onecolumn
\section{Conclusion: the useful convergence}
\label{sec:conclusion}
% =======================================================================

The useful convergence is not that Wolfram and the τ-framework say the
same thing. They do not. Wolfram's essay explores finite computational
microcosms and finds wildness: runtime outliers, isolate machines,
nondeterministic speedups, everything-machine reach, and irreducibility
pressure. The τ computation spine specifies an internal criterion under
which search becomes construction: finite interface width, structured
witness addresses, CRT decomposition, and Product-Meet Collapse.

The bridge is methodological. Both approaches reject the idea that
computation can be understood only from a fully abstract asymptotic
distance. Both insist that rule space, representation, observer access,
and infinity matter. Both make visible the hidden architecture under the
word ``search.''

But the difference is just as important. Wolfram shows how hard search can
look when finite rule universes are explored from below. The τ-framework
claims that search becomes structure only inside a specific admissible
address grammar. The note's title is therefore not a universal slogan. It
is a conditional research program:

\begin{quote}\small
Search becomes structure when the witness space has earned enough finite
address architecture for construction to replace blind enumeration.
\end{quote}

That is already a strong claim. It is strong enough to motivate a research
note, strong enough to compare with Wolfram's ruliological microcosms, and
bounded enough to send into public discussion without pretending that the
orthodox P-vs-NP problem has been settled.
This is an internal equality for τ-admissible witness problems. It does
not settle orthodox P vs NP over unrestricted Turing-machine / ZFC
encodings.

% =======================================================================
\section*{Acknowledgements}
\label{sec:ack}
% =======================================================================

This note was prepared as part of the Panta Rhei Research Program
research-note series.

% ---------------------------------------------------------------------------
\bibliographystyle{plain}
\bibliography{references}

\end{document}
